\section{Recorte, varianza e inferencia}\label{sec:inferencia}

\subsection{Recorte de pesos}\label{sec:recorte}

Recortar los pesos cambia el estimando. Para $\tau>0$, sean $w_\tau=\min(w,\tau)$ y $\varrho_\tau(t)=\min(1,\tau/t)$ para $t>0$. Con los pesos verdaderos, el estimador autonormalizado recortado converge a
\[
\theta_\tau=\frac{\E[w_\tau(X)\,Y\mid D=0]}{\E[w_\tau(X)\mid D=0]}.
\]

\begin{proposicion}[Sesgo del recorte]\label{prop:recorte}
Bajo (A1)--(A4):
\begin{enumerate}
  \item $\displaystyle\theta_\tau-\theta=\frac{\Cov_P\bigl(m_0(X),\varrho_\tau(w(X))\bigr)}{\E_P[\varrho_\tau(w(X))]}$.
  \item Si $h(t)=\E_P[m_0(X)\mid w(X)=t]$ admite una versión no creciente, $\theta_\tau\ge\theta$; si admite una no decreciente, $\theta_\tau\le\theta$.
  \item Si $Y\ge0$, la versión sin normalizar está siempre por debajo: $\E[w_\tau(X)\,Y\mid D=0]\le\theta$.
\end{enumerate}
\end{proposicion}
\begin{proof}
Como $P(\{w=0\})=\int_{\{w=0\}}w\dd Q=0$, sobre $\{w>0\}$ se tiene $w_\tau=\varrho_\tau(w)\,w$. Por la torre y el cambio de medida, $\E[w_\tau(X)Y\mid D=0]=\int w_\tau m_0\dd Q=\E_P[\varrho_\tau(w)\,m_0]$ y $\E[w_\tau(X)\mid D=0]=\E_P[\varrho_\tau(w)]\in(0,1]$; restando $\theta=\E_P[m_0]$ se obtiene (1). Para (2), la torre da $\Cov_P(m_0,\varrho_\tau(w))=\Cov_P(h(T),\varrho_\tau(T))$ con $T=w(X)$. La función $\varrho_\tau$ es no creciente; si $h$ también lo es, $(h(T)-h(T'))(\varrho_\tau(T)-\varrho_\tau(T'))\ge0$ para una copia independiente $T'$, y su esperanza es $2\Cov_P(h(T),\varrho_\tau(T))$ (desigualdad de asociación de Chebyshev); el caso no decreciente es simétrico. Para (3), $\varrho_\tau\le1$ y $m_0\ge0$ dan $\E_P[\varrho_\tau(w)m_0]\le\theta$.
\end{proof}

\paragraph{Interpretación.} $w(x)$ grande significa que el perfil $x$ es mucho más frecuente en el grupo objetivo que en el canal donante. No hay un sentido por defecto: depende de si, dentro del grupo objetivo, los perfiles más típicos del canal premium responden más o menos en el canal básico. Si responden menos, $h$ es decreciente y el recorte sobreestima $\theta$ y, por tanto, infraestima la pérdida; en nuestra simulación responden más (sección~\ref{sec:res-recorte}). La condición trata de $h$ como función del peso, no de $X$, y es aproximadamente comprobable, porque $m_0$ está identificada en el soporte.

\paragraph{Ejemplo~\ref{ej:juguete} (recorte).} Con $\tau=2$, los pesos $(0,2;\,1;\,3)$ pasan a $(0,2;\,1;\,2)$ y $\theta_\tau=1,2\,\%/0,8=1,5\,\%$, frente a $\theta=1,3\,\%$. La fórmula da lo mismo: $\varrho_\tau=(1;\,1;\,2/3)$, $\E_P[\varrho_\tau]=0,8$ y la covarianza vale $0,16$ puntos, así que el sesgo es $0,2$ puntos. La versión sin normalizar da $1,2\,\%$, por debajo de $\theta$.

\begin{corolario}[Recorte con el estimador doblemente robusto]\label{cor:drrecorte}
Por la Proposición~\ref{prop:dr} con $\bar w=w_\tau$, $\psi(w_\tau,\bar m)-\theta=-\int_{\{w>\tau\}}(w-\tau)(m_0-\bar m)\dd Q$.
\end{corolario}
El corolario se refiere a la forma no normalizada $\psi$: el recorte solo sesga en la medida en que el modelo de resultado se equivoca en la región recortada. Para la forma normalizada que usamos, el sesgo es $\E_Q[w_\tau(m_0-\bar m)]/\E_Q[w_\tau]-\E_P[m_0-\bar m]$, que se anula si $\bar m=m_0$ en todo el soporte, pero no necesariamente si el modelo solo acierta en la región recortada. Como la Proposición~\ref{prop:dr} vale para cualesquiera pesos, al recortar pesos estimados el sesgo depende de los dos errores, y se anula si $\bar m=m_0$.

\paragraph{Pesos estimados.} En la práctica se recorta $\hat w$, y el recorte elimina a la vez peso verdadero y ruido de estimación, así que el signo de la Proposición~\ref{prop:recorte} no está garantizado. En una nota complementaria en preparación mostramos que se conserva si el error de estimación es un ruido multiplicativo independiente del perfil con densidad log-cóncava. Antes de recortar, el efecto del recorte calculado con $\hat m$ en lugar de $Y$, o $\hat m$ por deciles de masa de $\hat w$, indica en qué sentido empujará. \citet{crump2009} recortan por la propensión y redefinen el estimando; \citet{ma2020} mantienen el estimando y caracterizan el sesgo del recorte por descarte en el estimador sin normalizar, del que la Proposición~\ref{prop:recorte}(3) es la versión para el recorte por tope; la aportación de la Proposición~\ref{prop:recorte} es el caso normalizado, con su fórmula exacta. \citet{liu2023} tratan el truncamiento como un cambio de estimando para el efecto sobre los tratados.

\subsection{Varianza y normalidad asintótica}

\paragraph{Función de influencia.} Sean $o=p/(1-p)$ y $\sigma^2(x)=\Var(Y\mid X=x,D=0)$. La función de influencia del estimador doblemente robusto, con la misma forma que la eficiente para la media contrafactual de los tratados \citep{hahn1998}, es
\begin{gather*}
\phi(O)=\frac1\pi\Bigl[Z\,(m_0(X)-\theta)+(1-Z)\,o(X)\,(Y-m_0(X))\Bigr],\\
\mathcal V=\E[\phi^2]=\frac{\Var_P(m_0(X))}{\pi}+\frac{\E_Q[w(X)^2\sigma^2(X)]}{1-\pi},
\end{gather*}
donde el término cruzado se anula porque $Z(1-Z)=0$ y el segundo sumando usa $o=\pi w/(1-\pi)$ (Teorema~\ref{teo:odds}). Esta forma cerrada corresponde a observaciones i.i.d.; con varias observaciones por unidad, véase el apéndice.

\begin{proposicion}[Cuándo la varianza es finita]\label{prop:varianza}
Si $Y$ está acotada y $0<\sigma_-^2\le\sigma^2\le\sigma_+^2$ $Q$-casi seguro, entonces $\mathcal V<\infty$ si y solo si $\E_Q[w^2]<\infty$, es decir, si y solo si $\chi^2(P\|Q)=\E_Q[w^2]-1<\infty$. La continuidad absoluta $P\ll Q$ no implica esa condición.
\end{proposicion}
\begin{proof}
El primer sumando de $\mathcal V$ es finito porque $m_0$ está acotada; el segundo está entre $\sigma_-^2\,\E_Q[w^2]/(1-\pi)$ y $\sigma_+^2\,\E_Q[w^2]/(1-\pi)$. Para ver que $P\ll Q$ no basta, sea $Q$ uniforme en $(0,1)$ y $w(x)=x^{-1/2}/2$, que integra 1; entonces $P\ll Q$ pero $\E_Q[w^2]=\int_0^1(4x)^{-1}dx=\infty$.
\end{proof}

Así, (A3) basta para identificar $\theta$, pero no para estimarla a ritmo $n^{-1/2}$: con la varianza condicional acotada inferiormente, hace falta además $\chi^2(P\|Q)<\infty$, y la teoría estándar exige algo más fuerte, pesos acotados.

\paragraph{Tamaño muestral efectivo.} El $\neff$ de Kish (sección~\ref{sec:estimadores}) cumple con los pesos verdaderos $\neff/n_0\to1/\E_Q[w^2]=1/(1+\chi^2(P\|Q))$, así que estima la divergencia $\chi^2$; con pesos estimados puede moverse en los dos sentidos (el ruido de $\hat w$ lo reduce y el encogimiento de un clasificador regularizado lo aumenta), y en nuestra simulación queda por debajo del teórico salvo con solapamiento bajo, donde el propio valor teórico es impreciso (sección~\ref{sec:diseno}). En muestra finita siempre es finito, incluso cuando $\chi^2=\infty$, y entonces tiende a cero muy despacio y con mucho ruido: en el contraejemplo de la Proposición~\ref{prop:varianza}, la mediana sobre réplicas pasa de 0,45 con $n_0=10^3$ a unos 0,22 con $n_0=10^7$, y con $n_0=10^3$ los percentiles 1 y 99 entre réplicas son 0,05 y 0,65. Un $\neff$ razonable no garantiza que $\chi^2$ sea finita; conviene examinar además la estabilidad de los cuantiles altos de $\hat w$.

\begin{teorema}[Normalidad asintótica del estimador doblemente robusto]\label{teo:normalidad}
Supongamos (A1)--(A4), $Y$ acotada y además:
\begin{enumerate}
  \item \emph{Solapamiento fuerte:} $p(X)\le1-\eta$ $Q$-casi seguro para algún $\eta>0$, es decir, $w$ acotada.
  \item \emph{Cross-fitting} por unidad en un número fijo $K$ de bloques de tamaño proporcional a $n$, con $\hat p\le1-\eta$ y $\hat m$ acotada por las mismas cotas que $Y$.
  \item \emph{Tasas:} $r_o=\max_k\|\hat o_k-o\|$ y $r_m=\max_k\|\hat m_k-m_0\|$ en $L^2(Q)$ tienden a 0 en probabilidad, y $r_o\,r_m=o_P(n^{-1/2})$.
\end{enumerate}
Entonces $\sqrt n(\hat\theta_{\mathrm{DR}}-\theta)\to N(0,\mathcal V)$ en distribución, y la media muestral de $\hat\phi(O_i)^2$ estima $\mathcal V$ de forma consistente.
\end{teorema}
La prueba, en el apéndice~\ref{app:prueba}, es una adaptación del Teorema 3.1 de \citet{chernozhukov2018} y de su ejemplo del efecto sobre los tratados: por la Proposición~\ref{prop:dr}, el sesgo es bilineal en los errores de los parámetros auxiliares (ortogonalidad de Neyman), el \emph{cross-fitting} controla el término de proceso empírico sin condiciones de complejidad, y el denominador se trata con el método delta.

\paragraph{Más allá del solapamiento fuerte.} La condición 1 es más fuerte que (A3) y que $\chi^2<\infty$. Si los pesos no están acotados pero $\chi^2<\infty$, $\mathcal V$ es finita, aunque la prueba estándar no cubre el caso. Si $\chi^2=\infty$ y la varianza condicional está acotada inferiormente, ningún estimador regular converge a ritmo $\sqrt n$: es el caso de identificación irregular de \citet{khan2010}. \citet{ma2020} muestran que entonces el límite puede ser no normal y proponen inferencia por submuestreo, y \citet{ma2023} dan inferencia adaptativa para estimadores doblemente robustos con recorte y corrección de sesgo. Nuestra simulación (sección~\ref{sec:diseno}) está en ese régimen: su índice de canal, lineal en covariables normales, no está acotado dentro de la región objetivo, así que los pesos no están acotados aunque $\chi^2$ sea finita, y la implementación recorta $\hat p$ a $[10^{-6},1-10^{-6}]$ en lugar de a $1-\eta$; sus coberturas son evidencia fuera de las condiciones del Teorema~\ref{teo:normalidad}.

\paragraph{Intervalos.} Usamos $\hat\theta_{\mathrm{DR}}\pm1,96\,(\hat{\mathcal V}/n)^{1/2}$. El Teorema~\ref{teo:normalidad} no cubre el recorte de pesos ni el bootstrap, que evaluamos en la simulación (secciones~\ref{sec:res-recorte} y~\ref{sec:res-bootstrap}). Si el grupo objetivo se define con un umbral estimado $\hat c$, su incertidumbre añade un término a $\mathcal V$ que no incluimos: tratamos la regla de selección como fija.
